\section*{Розділ IX. Прикінцеві положення}
\addcontentsline{toc}{section}{Розділ IX. Прикінцеві положення}

\subsection*{9.1. Порядок внесення змін та доповнень}
\addcontentsline{toc}{subsection}{9.1. Порядок внесення змін та доповнень}
    9.1.1. Зміни та доповнення до цього Положення вносяться Конференцією студентів Університету (КСУ) більшістю не менше двох третин голосів від загального складу делегатів КСУ.

    9.1.2. Право ініціювати внесення змін та доповнень до цього Положення мають:

        \begin{enumerate}[label=\alph*)]
            \item Конференція студентів Університету (КСУ);
            \item Студентська рада Київського авіаційного інституту (СР КАІ);
            \item Студентська рада студмістечка (СР СМ);
            \item Студентська уповноважена делегація (СУД);
            \item Центральна виборча комісія студентів (ЦВКс).
        \end{enumerate}

    9.1.3. Проєкт змін та доповнень до цього Положення підлягає попередньому оприлюдненню та обговоренню серед студентів у порядку, визначеному Регламентом КСУ.

\subsection*{9.2. Набуття чинності}
\addcontentsline{toc}{subsection}{9.2. Набуття чинності}
    9.2.1. Це Положення набуває чинності з дати його затвердження Конференцією студентів Університету та офіційного оприлюднення на інформаційних ресурсах органів студентського самоврядування Університету.

    9.2.2. З моменту набуття чинності цим Положенням усі попередні нормативні акти з питань організації виборів до органів студентського самоврядування, що суперечать цьому Положенню, втрачають чинність.

\subsection*{9.3. Тлумачення Положення}
\addcontentsline{toc}{subsection}{9.3. Тлумачення Положення}
    9.3.1. Офіційне тлумачення положень цього документа надає Студентська уповноважена делегація (СУД) за зверненням ЦВКс, органів студентського самоврядування або студентів Університету.

    9.3.2. Роз'яснення СУД щодо застосування норм цього Положення є обов'язковими для всіх органів студентського самоврядування та учасників виборчого процесу.

    9.3.3. У спірних питаннях застосування цього Положення ЦВКс керується принципами демократичності, прозорості та забезпечення виборчих прав студентів.

\subsection*{9.4. Перехідні положення}
\addcontentsline{toc}{subsection}{9.4. Перехідні положення}
    9.4.1. Перший склад ЦВКс після набуття чинності цим Положенням формується КСУ протягом одного місяця з дня затвердження цього Положення.

    9.4.2. Вибори до органів студентського самоврядування, призначені до набуття чинності цим Положенням, проводяться за раніше встановленими процедурами до їх завершення.

    9.4.3. Нові вибори, призначені після набуття чинності цим Положенням, проводяться виключно за процедурами, встановленими цим Положенням.

    9.4.4. ЦВКс протягом трьох місяців з дня початку своєї діяльності розробляє та подає на затвердження КСУ технічні регламенти проведення різних типів виборів відповідно до цього Положення.

\subsection*{9.5. Оскарження рішень ЦВКс}
\addcontentsline{toc}{subsection}{9.5. Оскарження рішень ЦВКс}
    9.5.1. Рішення ЦВКс можуть бути оскаржені до КСУ протягом 5 робочих днів з дня їх прийняття або оприлюднення.

    9.5.2. Право на оскарження мають:

        \begin{enumerate}[label=\alph*)]
            \item Кандидати, щодо яких прийнято рішення;
            \item Органи студентського самоврядування;
            \item Студентські організації;
            \item Групи студентів у складі не менше 15 осіб.
        \end{enumerate}

    9.5.3. Скарга на рішення ЦВКс повинна містити:

        \begin{enumerate}[label=\alph*)]
            \item Відомості про скаржника;
            \item Рішення ЦВКс, що оскаржується;
            \item Обґрунтування незаконності або необґрунтованості рішення;
            \item Вимоги скаржника;
            \item Докази, що підтверджують доводи скарги.
        \end{enumerate}

    9.5.4. КСУ розглядає скарги на рішення ЦВКс протягом 10 робочих днів з дня їх надходження та може:

        \begin{enumerate}[label=\alph*)]
            \item Залишити скаргу без задоволення;
            \item Скасувати рішення ЦВКс повністю або частково;
            \item Зобов'язати ЦВКс прийняти нове рішення з урахуванням зауважень КСУ.
        \end{enumerate}

\subsection*{9.6. Реорганізація та ліквідація ЦВКс}
\addcontentsline{toc}{subsection}{9.6. Реорганізація та ліквідація ЦВКс}
    9.6.1. Реорганізація або ліквідація ЦВКс можлива лише за рішенням КСУ у випадках:

        \begin{enumerate}[label=\alph*)]
            \item Систематичного невиконання ЦВКс своїх функцій;
            \item Грубого порушення принципів виборчого права;
            \item Реорганізації системи органів студентського самоврядування Університету.
        \end{enumerate}

    9.6.2. У разі ліквідації ЦВКс всі її документи та матеріали передаються на зберігання СУД.

    9.6.3. При реорганізації системи органів студентського самоврядування питання правонаступництва ЦВКс вирішується рішенням КСУ.

\subsection*{9.7. Заключні положення}
\addcontentsline{toc}{subsection}{9.7. Заключні положення}
    9.7.1. Це Положення є невід'ємною частиною нормативно-правової бази органів студентського самоврядування Київського авіаційного інституту.

    9.7.2. Положення, не врегульовані цим документом, вирішуються відповідно до Положення про ОСС, Статуту Університету та законодавства України.

    9.7.3. ЦВКс сприяє розвитку демократичної культури та громадянської активності студентів Університету через забезпечення прозорих та справедливих виборчих процесів.